% ══════════════════════════════════════════════════════════════
% Chapter 22: Docker Compose Deployment
% ══════════════════════════════════════════════════════════════
\chapter{Deployment: Docker Compose}
\label{ch:docker-deployment}

\begin{learnbox}
\begin{itemize}
  \item Deploy a multi-node blockchain network using Docker Compose
  \item Configure container networking, port mapping, and volume management
  \item Scale services and manage deployment topologies for development and testing
  \item Troubleshoot common Docker Compose deployment issues
  \item Understand container naming, instance identification, and execution flow
  \item Implement sequential startup coordination for reliable node formation
\end{itemize}
\end{learnbox}

\lettrine{U}{p} to this point in the book, we have been running a single node from \code{cargo run}. But Bitcoin is a distributed system\textemdash you need multiple nodes to test peer discovery, \index{block propagation}block propagation, and chain synchronization. \index{Docker}Docker Compose lets you spin up a configurable number of interconnected nodes with a single command, each with its own storage volume, while sharing a common network. This chapter bridges the gap between ``the code compiles'' and ``the system actually works as a network.''

\begin{infobox}{Scope}
This chapter covers local development and small-scale Docker Compose deployments. For production-grade orchestration with autoscaling and high availability, see Chapter~\ref{ch:k8s-deployment} (Kubernetes).
\end{infobox}

% ──────────────────────────────────────────────────────────────
\section{Overview and Quick Start}
\label{sec:docker-overview}

\index{Docker}Docker Compose deployment is \textbf{code-driven}: the \code{\index{Docker}docker-compose.yml} file describes the container topology, while orchestration logic (instance numbering, port selection, sequential startup, hostname resolution) lives in the \textbf{entrypoint and helper scripts}.

\subsection{Container Topology}

Figure~\ref{fig:docker-topology} shows the typical multi-node layout:

\begin{figure}[htbp]
  \centering
  \begin{tikzpicture}[
    node distance = 1.0cm and 1.2cm,
    container/.style = {rectangle, rounded corners=6pt, draw=sectioncolor, fill=sectioncolor!5,
                        text width=2.8cm, align=center, minimum height=2.4cm, font=\sffamily\small},
    component/.style = {rectangle, rounded corners=3pt, draw=sectioncolor!60, fill=codebg,
                        text width=1.8cm, align=center, minimum height=0.55cm, font=\ttfamily\tiny},
    network arrow/.style = {<->,thick, color=brand},
  ]

  % Docker Network border
  \draw[dashed,color=sectioncolor!50,line width=2pt] (-5.5,-0.8) rectangle (5.5,3.2);
  \node[font=\small\sffamily\color{sectioncolor!70}, anchor=north west] at (-5.3,3.1) {Docker Network};

  % Node 1
  \node[container] (node1) at (-3.5, 1.5) {
    \textbf{Node 1} \\[6pt]
    \texttt{:8000}
  };
  \node[component] (web1) at (-3.5, 2.7) {Web API};
  \node[component] (p2p1) at (-3.5, 2.0) {P2P};
  \node[component] (store1) at (-3.5, 1.3) {Store};
  \node[font=\tiny\sffamily,color=sectioncolor] at (-3.5, 0.8) {\texttt{vol\_1}};

  % Node 2
  \node[container] (node2) at (0, 1.5) {
    \textbf{Node 2} \\[6pt]
    \texttt{:8001}
  };
  \node[component] (web2) at (0, 2.7) {Web API};
  \node[component] (p2p2) at (0, 2.0) {P2P};
  \node[component] (store2) at (0, 1.3) {Store};
  \node[font=\tiny\sffamily,color=sectioncolor] at (0, 0.8) {\texttt{vol\_2}};

  % Node 3
  \node[container] (node3) at (3.5, 1.5) {
    \textbf{Node 3} \\[6pt]
    \texttt{:8002}
  };
  \node[component] (web3) at (3.5, 2.7) {Web API};
  \node[component] (p2p3) at (3.5, 2.0) {P2P};
  \node[component] (store3) at (3.5, 1.3) {Store};
  \node[font=\tiny\sffamily,color=sectioncolor] at (3.5, 0.8) {\texttt{vol\_3}};

  % Network connections between nodes
  \draw[network arrow] (p2p1.east) -- (p2p2.west);
  \draw[network arrow] (p2p2.east) -- (p2p3.west);

  % Host mapping
  \draw[dashed,color=codeframe,line width=1pt] (-5.5,-1.2) -- (5.5,-1.2);
  \node[font=\small\sffamily\color{codeframe}] at (-3.5,-1.5) {localhost:8000};
  \node[font=\small\sffamily\color{codeframe}] at (0,-1.5) {localhost:8001};
  \node[font=\small\sffamily\color{codeframe}] at (3.5,-1.5) {localhost:8002};
  \node[font=\small\sffamily\color{codeframe}] at (-5.3,-1.8) {\textbf{Host}};

  \end{tikzpicture}
  \caption{Docker Compose Deployment Topology: Multi-Node Network with Storage Volumes}
  \label{fig:docker-topology}
\end{figure}

\subsection{Quick Start: Single Miner and Webserver}

Before starting, set up wallet addresses. Use either address pool (recommended) or direct assignment:

\begin{listing}[htbp]
\caption{Quick Start: Address Pool Method}
\label{lst:ch22-quickstart-pool}
\begin{minted}[fontsize=\footnotesize]{bash}
cd ci/docker-compose/configs

# Set up address pool
export WALLET_ADDRESS_POOL="addr1,addr2"

# Start services
docker compose up -d

# Result:
# - redis (rate limiting backend)
# - miner_1 (P2P port 2001, uses addr1)
# - webserver_1 (Web 8080, P2P 2101, uses addr2)
\end{minted}
\end{listing}

Alternative: direct assignment:

\begin{listing}[htbp]
\caption{Quick Start: Direct Assignment}
\label{lst:ch22-quickstart-direct}
\begin{minted}[fontsize=\footnotesize]{bash}
cd ci/docker-compose/configs
export NODE_MINING_ADDRESS=3npBNyKSEwhCQWTXHFjwR8Rb66kjq6khfZSdmLPm8Gde9XoTwW
docker compose up -d
\end{minted}
\end{listing}

Access webserver at \texttt{http://localhost:8080}.

\subsection{Scaling to Multiple Instances}

Once comfortable with basic setup, scale to multiple instances:

\begin{listing}[htbp]
\caption{Scale to Multiple Miners and Webservers}
\label{lst:ch22-scale}
\begin{minted}[fontsize=\footnotesize]{bash}
cd ci/docker-compose/configs

# Set up address pool (at least 5 addresses)
export WALLET_ADDRESS_POOL="addr1,addr2,addr3,addr4,addr5"

# Scale to 3 miners and 2 webservers
./docker-compose.scale.sh 3 2

# Results:
# - miner_1 -> addr1 (port 2001)
# - miner_2 -> addr2 (port 2002)
# - miner_3 -> addr3 (port 2003)
# - webserver_1 -> addr4 (port 8080, P2P 2101)
# - webserver_2 -> addr5 (port 8081, P2P 2102)
\end{minted}
\end{listing}

\begin{notebox}
Always use the helper script \code{./docker-compose.scale.sh} to ensure \textbf{ALL instances} have ports accessible externally. Using \code{docker compose --scale} directly only maps ports for the first instance.
\end{notebox}

% ──────────────────────────────────────────────────────────────
\section{Architecture and Execution Flow}
\label{sec:docker-architecture}

\subsection{Container Naming and Instance Identification}

Docker Compose names containers using the pattern:
\begin{equation}
\text{<project>\_<service>\_<instance>}
\end{equation}

\textbf{Examples}:
\begin{itemize}
  \item \code{blockchain\_miner\_1} (first miner instance)
  \item \code{blockchain\_miner\_2} (second miner instance)
  \item \code{blockchain\_webserver\_1} (first webserver instance)
\end{itemize}

The entrypoint script extracts container identity from Docker's \code{HOSTNAME} environment variable and normalizes it into \code{CONTAINER\_NAME}. It then extracts the instance number using a regex pattern:

\begin{listing}[htbp]
\caption{Instance Number Extraction}
\label{lst:ch22-instance-extract}
\begin{minted}[fontsize=\footnotesize]{bash}
CONTAINER_NAME="${HOSTNAME:-}"  # e.g. "blockchain_miner_1"

# Extract instance number from container name
if [[ "${CONTAINER_NAME}" =~ _([0-23]+)$ ]]; then
    INSTANCE_NUMBER="${BASH_REMATCH[1]}"  # e.g. 1
fi
\end{minted}
\end{listing}

\subsection{Port Calculation}

Based on instance number and service type, ports are calculated:

\begin{table}[htbp]
\caption{Port Calculation Formula}
\label{tab:port-calc}
\centering\small
\begin{tabularx}{0.85\textwidth}{lllX}
\toprule
\textbf{Service} & \textbf{Instance} & \textbf{Formula} & \textbf{Port} \\
\midrule
Miner & 1 & \(2001 + 1 - 1\) & 2001 \\
Miner & 2 & \(2001 + 2 - 1\) & 2002 \\
Miner & 3 & \(2001 + 3 - 1\) & 2003 \\
\midrule
Webserver (Web) & 1 & \(8080 + 1 - 1\) & 8080 \\
Webserver (Web) & 2 & \(8080 + 2 - 1\) & 8081 \\
Webserver (P2P) & 1 & \(2101 + 1 - 1\) & 2101 \\
Webserver (P2P) & 2 & \(2101 + 2 - 1\) & 2102 \\
\bottomrule
\end{tabularx}
\end{table}

\subsection{Data Directory Isolation}

Each instance uses an isolated data directory:

\begin{listing}[htbp]
\caption{Data Directory Naming}
\label{lst:ch22-data-dir}
\begin{minted}[fontsize=\footnotesize]{bash}
# Per-instance data directories
DATA_DIR="data${INSTANCE_NUMBER}"
BLOCKS_TREE="blocks${INSTANCE_NUMBER}"

# Examples:
# Instance 1: data1/, blocks1/
# Instance 2: data2/, blocks2/
# Instance 3: data3/, blocks3/
\end{minted}
\end{listing}

Data is persisted in Docker volumes:
\begin{itemize}
  \item Miners: \code{miner-data} and \code{miner-wallets}
  \item Webservers: \code{webserver-data} and \code{webserver-wallets}
\end{itemize}

\subsection{Execution Timeline}

The startup flow progresses through these phases:

\begin{listing}[htbp]
\caption{Docker Compose Startup Order}
\label{lst:ch22-exec-order}
\begin{minted}[fontsize=\footnotesize]{text}
1. docker-compose.yml (Docker Compose reads configuration)
   v
2. Dockerfile (builds the blockchain binary into an image)
   v
3. docker-entrypoint.sh (Container startup script)
   +- wait-for-node.sh (if sequential startup requires waiting)
   v
4. /app/blockchain startnode ... (Rust binary entry point)
\end{minted}
\end{listing}

% ──────────────────────────────────────────────────────────────
\section{Deployment Topology and Networking}
\label{sec:docker-topology}

\subsection{Miner Connection Chain}

When scaling miners, they form a sequential chain:

\begin{listing}[htbp]
\caption{Miner Network Topology}
\label{lst:ch22-miner-chain}
\begin{minted}[fontsize=\footnotesize]{text}
Miner 1: NODE_CONNECT_NODES=local (seed node)
  v
Miner 2: NODE_CONNECT_NODES=miner_1:2001
  v
Miner 3: NODE_CONNECT_NODES=miner_2:2002
  v
Miner 4: NODE_CONNECT_NODES=miner_3:2003
\end{minted}
\end{listing}

\subsection{Webserver Connections}

\textbf{Important}: Webservers \textbf{always} connect to miners, never to other webservers.

\begin{listing}[H]
\caption{Webserver Network Topology}
\label{lst:ch22-webserver-topology}
\begin{minted}[fontsize=\footnotesize]{text}
Webserver 1: NODE_CONNECT_NODES=miner_1:2001
Webserver 2: NODE_CONNECT_NODES=miner_1:2001 (first miner, not webserver_1)
Webserver 3: NODE_CONNECT_NODES=miner_1:2001 (first miner, not webserver_2)
\end{minted}
\end{listing}

\subsection{Network Topologies}

\subsubsection{Single Miner + Multiple Webservers (Star)}

All webservers connect to the single miner:

\begin{listing}[htbp]
\caption{Star Topology: One Miner, Multiple Webservers}
\label{lst:ch22-star-topology}
\begin{minted}[fontsize=\footnotesize]{bash}
./docker-compose.scale.sh 1 3
\end{minted}
\end{listing}

\subsubsection{Multiple Miners (Chain)}

Miners form a linear chain:

\begin{listing}[htbp]
\caption{Chain Topology: Multiple Miners}
\label{lst:ch22-chain-topology}
\begin{minted}[fontsize=\footnotesize]{bash}
./docker-compose.scale.sh 3 1
\end{minted}
\end{listing}

\subsubsection{Multiple Miners + Multiple Webservers (Hybrid)}

Miners form a chain; webservers star around the first miner:

\begin{listing}[htbp]
\caption{Hybrid Topology: Multiple Miners and Webservers}
\label{lst:ch22-hybrid-topology}
\begin{minted}[fontsize=\footnotesize]{bash}
./docker-compose.scale.sh 3 2
\end{minted}
\end{listing}

\subsection{Sequential Startup Coordination}

When \code{SEQUENTIAL\_STARTUP=yes} (default), nodes start in sequence:

\begin{enumerate}
  \item Instance 1 starts immediately (no wait needed)
  \item Instance 2 waits for Instance 1 to be ready, then connects to it
  \item Instance 3 waits for Instance 2 to be ready, then connects to it
  \item And so on\ldots
\end{enumerate}

The \code{wait-for-node.sh} script handles this coordination. For miners, it waits for the previous miner. For webservers, it always waits for miners.

\begin{listing}[htbp]
\caption{Sequential Startup Configuration}
\label{lst:ch22-sequential-config}
\begin{minted}[fontsize=\footnotesize]{bash}
# Enable sequential startup (default)
export SEQUENTIAL_STARTUP=yes
./docker-compose.scale.sh 3 2

# Or disable for simultaneous startup
export SEQUENTIAL_STARTUP=no
./docker-compose.scale.sh 3 2
\end{minted}
\end{listing}

% ──────────────────────────────────────────────────────────────
\section{Port Mapping and External Access}
\label{sec:docker-ports}

\subsection{The Port Mapping Challenge}

Docker Compose only maps ports for the first instance when using \code{--scale}. Additional instances created by scaling do NOT get ports mapped to the host.

\begin{table}[htbp]
\caption{Port Mapping Limitation}
\label{tab:port-mapping-issue}
\centering\small
\begin{tabularx}{0.85\textwidth}{llll}
\toprule
\textbf{Instance} & \textbf{Internal} & \textbf{External} & \textbf{Accessible?} \\
\midrule
miner\_1 & 2001 & 2001 & Yes \\
miner\_2 & 2002 & --- & No \\
miner\_3 & 2003 & --- & No \\
\bottomrule
\end{tabularx}
\end{table}

\subsection{Solution: Use Helper Script}

The \code{./docker-compose.scale.sh} script automatically generates port mappings for all instances:

\begin{listing}[htbp]
\caption{Using the Helper Script for Scaling}
\label{lst:ch22-helper-script}
\begin{minted}[fontsize=\footnotesize]{bash}
cd ci/docker-compose/configs
./docker-compose.scale.sh 3 2

# This script:
# 1. Generates docker-compose.override.yml with port mappings
# 2. Scales services with all ports accessible externally
# 3. All instances have ports mapped to host
\end{minted}
\end{listing}

\needspace{3in}
\subsection{Alternative: Manual Override File}

If you prefer to generate the override file manually:

\begin{listing}[htbp]
\caption{Manual Port Override Generation}
\label{lst:ch22-manual-override}
\begin{minted}[fontsize=\footnotesize]{bash}
cd ci/docker-compose/configs

# Step 1: Generate override file
./generate-compose-ports.sh 3 2

# Step 2: Start services
docker compose up -d --scale miner=3 --scale webserver=2
\end{minted}
\end{listing}

% ──────────────────────────────────────────────────────────────
\section{DNS Resolution Mechanism}
\label{sec:docker-dns}

\subsection{Understanding Docker Compose DNS}

Docker Compose creates an internal DNS system that resolves \textbf{service names} (defined in \code{docker-compose.yml}), but \textbf{NOT instance names}.

\begin{table}[htbp]
\caption{Docker Compose DNS Resolution}
\label{tab:docker-dns-resolution}
\centering\small
\begin{tabularx}{0.85\textwidth}{llll}
\toprule
\textbf{Hostname} & \textbf{Type} & \textbf{Resolves?} & \textbf{Result} \\
\midrule
\code{miner} & Service name & Yes & 172.19.0.2 \\
\code{miner\_1} & Instance name & No & Error \\
\code{miner\_2} & Instance name & No & Error \\
\bottomrule
\end{tabularx}
\end{table}

\subsection{The Conversion Mechanism}

The entrypoint script converts instance names to service names for DNS resolution:

\begin{listing}[htbp]
\caption{Instance Name to Service Name Conversion}
\label{lst:ch22-dns-conversion}
\begin{minted}[fontsize=\footnotesize]{bash}
# Input from wait script
PREV_ADDR="miner_1:2001"

# Initialize resolution address
RESOLVE_ADDR="${PREV_ADDR}"

# Pattern match: miner_N:PORT
if [[ "${PREV_ADDR}" =~ ^miner_([0-23]+): ]]; then
    # Extract port and convert to service name
    PORT_PART="${PREV_ADDR##*:}"          # Extract port
    RESOLVE_ADDR="miner:${PORT_PART}"     # Use service name
fi

# Now resolve using Docker DNS
PREV_ADDR_RESOLVED=$(resolve_hostname_to_ip "${RESOLVE_ADDR}")
# Result: PREV_ADDR_RESOLVED="172.19.0.2:2001"
\end{minted}
\end{listing}

\textbf{Why this works}: \code{miner} (the service name) resolves in Docker Compose DNS, but \code{miner\_1} (the instance name) does not.

% ──────────────────────────────────────────────────────────────
\section{Deployment Scenarios and Operations}
\label{sec:docker-scenarios}

\subsection{Scenario 1: Development Setup}

Single miner and webserver for local development:

\begin{listing}[htbp]
\caption{Single Miner + Webserver Deployment}
\label{lst:ch22-scenario-dev}
\begin{minted}[fontsize=\footnotesize]{bash}
cd ci/docker-compose/configs

export WALLET_ADDRESS_POOL="addr1,addr2"
docker compose up -d

# Accessible at:
# - Miner P2P: localhost:2001
# - Webserver Web: http://localhost:8080
# - Webserver P2P: localhost:2101
\end{minted}
\end{listing}

\subsection{Scenario 2: Multiple Mining Nodes}

Test distributed mining with three miners:

\begin{listing}[htbp]
\caption{Three Miners + One Webserver}
\label{lst:ch22-scenario-mining}
\begin{minted}[fontsize=\footnotesize]{bash}
cd ci/docker-compose/configs

export WALLET_ADDRESS_POOL="addr1,addr2,addr3,addr4"
./docker-compose.scale.sh 3 1

# Topology:
# miner_1 <- miner_2 <- miner_3
# ^
# webserver_1 (accesses all blocks via miner network)
\end{minted}
\end{listing}

\subsection{Scenario 3: High Availability Web Access}

Multiple webservers for redundancy:

\begin{listing}[htbp]
\caption{One Miner + Three Webservers}
\label{lst:ch22-scenario-webservers}
\begin{minted}[fontsize=\footnotesize]{bash}
cd ci/docker-compose/configs

export WALLET_ADDRESS_POOL="addr1,addr2,addr3,addr4"
./docker-compose.scale.sh 1 3

# Result:
# - miner_1 (port 2001)
# - webserver_1 (port 8080)
# - webserver_2 (port 8081)
# - webserver_3 (port 8082)
# All webservers connect to miner_1
\end{minted}
\end{listing}

\subsection{Scenario 4: Production-like Configuration}

Multiple miners and webservers:

\begin{listing}[htbp]
\caption{Three Miners + Two Webservers}
\label{lst:ch22-scenario-production}
\begin{minted}[fontsize=\footnotesize]{bash}
cd ci/docker-compose/configs

export WALLET_ADDRESS_POOL="addr1,addr2,addr3,addr4,addr5"
./docker-compose.scale.sh 3 2
\end{minted}
\end{listing}

% ──────────────────────────────────────────────────────────────
\needspace{2.5in}
\section{Managing Script Changes}
\label{sec:docker-script-changes}

\subsection{Important: Rebuild Docker Image}

The \code{docker-entrypoint.sh} and \code{wait-for-node.sh} scripts are copied into the Docker image during build. To deploy changes, you \textbf{MUST rebuild} the Docker image---Docker will use cached layers otherwise.

\begin{listing}[htbp]
\caption{Deploying Script Changes}
\label{lst:ch22-rebuild}
\begin{minted}[fontsize=\footnotesize]{bash}
# CRITICAL: Use --no-cache for script changes
cd ci/docker-compose/configs

# Stop existing containers
docker compose down

# Rebuild with no cache
docker compose build --no-cache

# Start fresh
docker compose up -d
\end{minted}
\end{listing}

\begin{warnbox}[Caching Pitfall]
Using \code{docker compose build} alone may use cached layers. For script changes, always use \code{--no-cache}.
\end{warnbox}

% ──────────────────────────────────────────────────────────────
\section{Environment Variables and Configuration}
\label{sec:docker-config}

\subsection{Wallet Address Distribution}

\subsubsection{Method 1: Address Pool (Recommended)}

Each node automatically selects an address by index:

\begin{listing}[htbp]
\caption{Address Pool Configuration}
\label{lst:ch22-address-pool}
\begin{minted}[fontsize=\footnotesize]{bash}
# Create pool with at least one address per instance
export WALLET_ADDRESS_POOL="addr1,addr2,addr3,addr4,addr5"

# Start with automatic selection
docker compose up -d --scale miner=3 --scale webserver=2

# Result:
# miner_1 -> addr1
# miner_2 -> addr2
# miner_3 -> addr3
# webserver_1 -> addr4
# webserver_2 -> addr5
\end{minted}
\end{listing}

\subsubsection{Method 2: Direct Assignment}

Use the same address for all instances:

\begin{listing}[htbp]
\caption{Direct Address Assignment}
\label{lst:ch22-direct-address}
\begin{minted}[fontsize=\footnotesize]{bash}
export NODE_MINING_ADDRESS=3npBNyKSEwhCQWTXHFjwR8Rb66kjq6khfZSdmLPm8Gde9XoTwW
docker compose up -d --scale miner=3
\end{minted}
\end{listing}

\subsection{Rate Limiting Configuration}

The webserver loads \code{axum\_rate\_limiter} settings from \code{ci/docker-compose/configs/Settings.toml}:

\begin{listing}[htbp]
\caption{Rate Limiting Setup}
\label{lst:ch22-rate-limit}
\begin{minted}[fontsize=\footnotesize]{bash}
# Compose mounts Settings.toml -> /app/Settings.toml
# Environment variable: RL_SETTINGS_PATH=/app/Settings.toml

# To change limits:
# 1. Edit configs/Settings.toml
# 2. Restart webserver
docker compose restart webserver
\end{minted}
\end{listing}

% ──────────────────────────────────────────────────────────────
\section{Monitoring and Troubleshooting}
\label{sec:docker-troubleshooting}

\subsection{View Running Containers}

\begin{listing}[H]
\caption{Check Container Status}
\label{lst:ch22-status}
\begin{minted}[fontsize=\footnotesize]{bash}
docker compose ps
\end{minted}
\end{listing}

\subsection{View Logs}

\begin{listing}[H]
\caption{Monitor Container Logs}
\label{lst:ch22-logs}
\begin{minted}[fontsize=\footnotesize]{bash}
# All miner logs
docker compose logs -f miner

# All webserver logs
docker compose logs -f webserver

# Specific instance
docker compose logs -f miner_1
\end{minted}
\end{listing}

\subsection{Common Issues}

\subsubsection{Port Conflicts}

\begin{listing}[H]
\caption{Diagnose Port Conflicts}
\label{lst:ch22-port-conflicts}
\begin{minted}[fontsize=\footnotesize]{bash}
docker compose ps
netstat -tulpn | grep -E '2001|8080|2101'
\end{minted}
\end{listing}

\subsubsection{Data Persistence}

\begin{listing}[H]
\caption{Check Data Persistence}
\label{lst:ch22-data-check}
\begin{minted}[fontsize=\footnotesize]{bash}
# List volumes
docker volume ls | grep blockchain

# Inspect volume
docker volume inspect blockchain_miner-data

# Remove volumes (WARNING: deletes data)
docker compose down -v
\end{minted}
\end{listing}

\subsubsection{Nodes Not Connecting}

\begin{listing}[H]
\caption{Troubleshoot Node Connectivity}
\label{lst:ch22-connectivity}
\begin{minted}[fontsize=\footnotesize]{bash}
# Check if target node is ready
docker compose exec miner_2 nc -zv miner_1 2001

# Check connection environment variable
docker compose exec miner_2 env | grep NODE_CONNECT_NODES

# View connection errors in logs
docker compose logs miner_2 | grep -i connect
\end{minted}
\end{listing}

% ──────────────────────────────────────────────────────────────
\section{Accessing the Webserver}
\label{sec:docker-webserver-access}

\subsection{Health Checks}

\begin{listing}[H]
\caption{Webserver Health Endpoints}
\label{lst:ch22-health}
\begin{minted}[fontsize=\footnotesize]{bash}
# Liveness probe
curl -f http://localhost:8080/api/health/live

# Readiness probe
curl -f http://localhost:8080/api/health/ready
\end{minted}
\end{listing}

\subsection{Swagger UI}

Access interactive API documentation at:
\begin{equation}
\text{\texttt{http://localhost:8080/swagger-ui/}}
\end{equation}

\subsection{API Authentication}

All API endpoints require the \code{X-API-Key} header. Default keys:
\begin{itemize}
  \item Admin API: \code{admin-secret}
  \item Wallet API: \code{wallet-secret}
\end{itemize}

\begin{listing}[htbp]
\caption{API Call Examples}
\label{lst:ch22-api-examples}
\begin{minted}[fontsize=\footnotesize]{bash}
# Get blockchain info
curl -H "X-API-Key: admin-secret" \
  http://localhost:8080/api/admin/blockchain

# Get latest blocks
curl -H "X-API-Key: admin-secret" \
  http://localhost:8080/api/admin/blockchain/blocks/latest

# Create wallet
curl -X POST \
  -H "X-API-Key: wallet-secret" \
  -H "Content-Type: application/json" \
  http://localhost:8080/api/wallet/wallet
\end{minted}
\end{listing}

% ──────────────────────────────────────────────────────────────
\section{Cleanup}
\label{sec:docker-cleanup}

\subsection{Stop Services}

\begin{listing}[H]
\caption{Stopping Docker Compose Services}
\label{lst:ch22-cleanup}
\begin{minted}[fontsize=\footnotesize]{bash}
# Stop all services (keeps volumes)
docker compose down

# Stop and remove volumes (deletes all data)
docker compose down -v

# View volumes before deletion
docker volume ls | grep blockchain
\end{minted}
\end{listing}

% ──────────────────────────────────────────────────────────────
\begin{chaptersummary}
\item We deployed a multi-node blockchain network using Docker Compose with configurable topologies.
\item We configured container networking, port mapping via helper scripts, and volume management for data persistence.
\item We understood instance identification, port calculation, and data directory isolation.
\item We implemented sequential startup coordination to ensure reliable node formation.
\item We managed DNS resolution by converting instance names to service names for Docker's internal DNS.
\item We deployed script changes by rebuilding Docker images without cache.
\item We monitored and troubleshot common deployment issues with logging and health checks.
\end{chaptersummary}

---
